\documentclass[aspectratio=169]{beamer}

\usepackage[utf8]{inputenc}
\usepackage[T1]{fontenc}
\usepackage{amsmath,amssymb,amsfonts,bm,mathtools}
\usepackage{physics}
\usepackage{hyperref}
\usepackage{booktabs}

\usetheme{Madrid}

\title{Review of Chapters 1--2}
\author{MHJ inputs to Antons}
\date{April 27}

\begin{document}

\frame{\titlepage}

%================================================
\section{Overall assessment}
%================================================

\begin{frame}{Overall Assessment}
The first two chapters present a strong and technically ambitious MSc thesis project. They are well written and at the right level. 

\medskip

The work combines:
\begin{itemize}
    \item heat transfer,
    \item electromagnetics,
    \item finite element methods,
    \item physics-informed machine learning,
    \item graph neural networks,
    \item industrial induction heating applications.
\end{itemize}
\end{frame}

\begin{frame}{General Strengths}
As far as the first read your drafts of chapters 1 and 2 show:
\begin{itemize}
    \item good interdisciplinary scope,
    \item clear computational science framework,
    \item relevant industrial motivation,
    \item sensible modular decomposition into electromagnetic and thermal surrogate models,
    \item awareness of verification through analytical tests and manufactured solutions.
\end{itemize}

\medskip

The text is strongest when it connects physical modeling, numerical methods, and machine-learning architecture.
\end{frame}

\begin{frame}{Things that can be improved}
The main topics to improve are  not the scientific directions, but rather some details in the presentation.

\medskip

At present, chapters 1 and 1 sometimes read as:
\[
\text{technically correct summary}
\quad\text{rather than}\quad
\text{fully developed scientific argument}.
\]

\medskip

The most important improvement is to strengthen the narrative:
\[
\text{physics problem}
\rightarrow
\text{computational bottleneck}
\rightarrow
\text{surrogate-model strategy}
\rightarrow
\text{physics-informed graph-learning formulation}.
\]
\end{frame}

%================================================
\section{Chapter 1}
%================================================

\begin{frame}{Chapter 1: Theoretical Background}
Chapter 1 covers the correct topics:
\begin{itemize}
    \item heat equation,
    \item Maxwell equations and magnetoquasistatics,
    \item induction heating,
    \item finite element methods,
    \item physics-informed neural networks,
    \item graph neural networks.
\end{itemize}

\medskip

This is the right theoretical foundation for the thesis.

\medskip

However, the chapter would benefit from more mathematical depth and stronger transitions between topics.
\end{frame}

\begin{frame}{Chapter 1: Strengths}
Strengths:
\begin{itemize}
    \item The scope is appropriate.
    \item The physical motivation is clear.
    \item The connection between induction heating and coupled electromagnetic--thermal physics is well chosen.
    \item FEM is correctly introduced as the conventional numerical benchmark.
    \item PINNs and GNNs are relevant to the computational problem.
\end{itemize}
\end{frame}

\begin{frame}{Chapter 1: Main imporvements}
The chapter is occasionally too compact and catalog-like.

\medskip

It would be improved by explicitly explaining why each ingredient is needed.

\medskip

For example:
\[
\text{induction heating}
\Rightarrow
\text{Maxwell equations + heat equation}
\Rightarrow
\text{FEM benchmark}
\Rightarrow
\text{computational cost}
\Rightarrow
\text{physics-informed surrogate model}.
\]

\medskip

This would make the chapter read more like a thesis argument and less like a list of background topics.
\end{frame}

%================================================
\section{Equations for Chapter 1}
%================================================

\begin{frame}{Suggested Equation for Section 1.1.1: Heat Equation}
A central equation to include explicitly is the transient heat equation:
\[
\rho c_p \frac{\partial T}{\partial t}
=
\nabla\cdot(k\nabla T) + Q.
\]

Here:
\begin{itemize}
    \item \(T(\mathbf{x},t)\) is temperature,
    \item \(\rho\) is density,
    \item \(c_p\) is heat capacity,
    \item \(k\) is thermal conductivity,
    \item \(Q\) is the volumetric heat source.
\end{itemize}

\medskip

For induction heating, \(Q\) should later be connected to electromagnetic losses.
\end{frame}

\begin{frame}{Suggested Equation: Boundary and Initial Conditions}
The heat problem should be completed with boundary and initial data:
\[
T(\mathbf{x},0)=T_0(\mathbf{x}),
\]
\[
T = T_D \quad \text{on } \Gamma_D,
\]
\[
-k\nabla T\cdot \mathbf{n}=q_N \quad \text{on } \Gamma_N.
\]

For convective cooling, include Robin boundary conditions:
\[
-k\nabla T\cdot \mathbf{n}
=
h(T-T_\infty)
\quad \text{on } \Gamma_R.
\]

This makes the PDE problem mathematically complete.
\end{frame}

\begin{frame}{Suggested Equations for Section 1.1.2: Maxwell Equations}
The electromagnetic section should begin from Maxwell's equations:
\[
\nabla\cdot \mathbf{D}=\rho_e,
\]
\[
\nabla\cdot \mathbf{B}=0,
\]
\[
\nabla\times \mathbf{E}=-\frac{\partial \mathbf{B}}{\partial t},
\]
\[
\nabla\times \mathbf{H}=\mathbf{J}+\frac{\partial \mathbf{D}}{\partial t}.
\]

Then state the magnetoquasistatic approximation relevant for induction heating:
\[
\nabla\times \mathbf{H}\approx \mathbf{J}.
\]
\end{frame}

\begin{frame}{Suggested Equations: Constitutive Relations}
Include constitutive relations:
\[
\mathbf{B}=\mu \mathbf{H},
\qquad
\mathbf{J}=\sigma \mathbf{E}.
\]

For temperature-dependent materials:
\[
\sigma=\sigma(T),
\qquad
\mu=\mu(T),
\qquad
k=k(T).
\]

This is important because the thesis emphasizes temperature-dependent material properties and coupled electromagnetic-thermal modeling.
\end{frame}

\begin{frame}{Suggested Equations: Magnetic Vector Potential}
Use the magnetic vector potential:
\[
\mathbf{B}=\nabla\times \mathbf{A}.
\]

Then Faraday's law gives, in a suitable gauge,
\[
\mathbf{E}=-\frac{\partial \mathbf{A}}{\partial t}-\nabla \varphi.
\]

For frequency-domain harmonic fields:
\[
\mathbf{A}(\mathbf{x},t)=\Re\{\mathbf{A}(\mathbf{x})e^{i\omega t}\}.
\]

A common magnetoquasistatic equation is
\[
\nabla\times\left(\nu\nabla\times \mathbf{A}\right)
+
i\omega\sigma \mathbf{A}
=
\mathbf{J}_{\mathrm{src}},
\]
where \(\nu=1/\mu\).
\end{frame}

\begin{frame}{Suggested Equation for Section 1.1.3: Joule Heating}
The electromagnetic and thermal models should be connected through Joule heating.

For harmonic fields, a standard time-averaged source is
\[
Q_{\mathrm{Joule}}
=
\frac{1}{2}\sigma |\mathbf{E}|^2.
\]

Using the vector potential approximation,
\[
\mathbf{E}\approx -i\omega \mathbf{A},
\]
one obtains
\[
Q_{\mathrm{Joule}}
=
\frac{1}{2}\sigma \omega^2 |\mathbf{A}|^2.
\]

This equation is central for induction heating and should be highlighted.
\end{frame}

%================================================
\section{FEM equations}
%================================================

\begin{frame}{Suggested Equation for Section 1.2.1: Weak Heat Form}
The FEM section should include at least one full weak-form derivation.

Starting from
\[
\rho c_p \frac{\partial T}{\partial t}
-
\nabla\cdot(k\nabla T)
=
Q,
\]
multiply by a test function \(v\) and integrate:
\[
\int_\Omega \rho c_p v \frac{\partial T}{\partial t}\,d\Omega
-
\int_\Omega v\nabla\cdot(k\nabla T)\,d\Omega
=
\int_\Omega vQ\,d\Omega.
\]

Integrating by parts gives:
\[
\int_\Omega \rho c_p v \frac{\partial T}{\partial t}\,d\Omega
+
\int_\Omega k\nabla v\cdot \nabla T\,d\Omega
=
\int_\Omega vQ\,d\Omega
+
\int_{\partial\Omega} v k\nabla T\cdot \mathbf{n}\,dS.
\]
\end{frame}

\begin{frame}{Suggested Equation: FEM Matrix Form}
After finite-element discretization,
\[
T(\mathbf{x},t)\approx \sum_j T_j(t)N_j(\mathbf{x}),
\]
the semi-discrete system becomes
\[
M \dot{\mathbf{T}} + K\mathbf{T}=\mathbf{f},
\]
with
\[
M_{ij}=\int_\Omega \rho c_p N_iN_j\,d\Omega,
\]
\[
K_{ij}=\int_\Omega k\nabla N_i\cdot\nabla N_j\,d\Omega,
\]
\[
f_i=\int_\Omega N_i Q\,d\Omega.
\]

This gives an explicit bridge between FEM and graph-based learning on mesh nodes.
\end{frame}

\begin{frame}{Suggested Equation for Section 1.2.2: Time Discretization}
For backward Euler:
\[
M\frac{\mathbf{T}^{n+1}-\mathbf{T}^{n}}{\Delta t}
+
K\mathbf{T}^{n+1}
=
\mathbf{f}^{n+1}.
\]

Equivalently:
\[
\left(M+\Delta t K\right)\mathbf{T}^{n+1}
=
M\mathbf{T}^{n}
+
\Delta t\,\mathbf{f}^{n+1}.
\]

This is important because the thermal surrogate appears to learn a time-stepping map:
\[
\mathbf{T}^{n}\mapsto \mathbf{T}^{n+1}.
\]
\end{frame}

\begin{frame}{Suggested Equation for Section 1.2.3: Coupled EM--Thermal Loop}
A clear coupled model can be written as:
\[
\nabla\times\left(\nu(T)\nabla\times \mathbf{A}\right)
+
i\omega\sigma(T)\mathbf{A}
=
\mathbf{J}_{\mathrm{src}},
\]
\[
\rho c_p(T)\frac{\partial T}{\partial t}
=
\nabla\cdot(k(T)\nabla T)
+
\frac{1}{2}\sigma(T)\omega^2|\mathbf{A}|^2.
\]

This pair of equations summarizes the physical coupling targeted by the thesis.
\end{frame}

%================================================
\section{ML equations}
%================================================

\begin{frame}{Suggested Equation for Section 1.3: Neural Network Map}
Introduce the supervised-learning viewpoint:
\[
\hat{\mathbf{y}} = f_\theta(\mathbf{x}),
\]
with a data loss
\[
\mathcal{L}_{\mathrm{data}}(\theta)
=
\frac{1}{N}\sum_{i=1}^N
\norm{f_\theta(\mathbf{x}_i)-\mathbf{y}_i}^2.
\]

Then explain that physics-informed learning augments this with PDE residuals.
\end{frame}

\begin{frame}{Suggested Equation for Section 1.3.1: PINN Loss}
For a PDE
\[
\mathcal{N}[u](\mathbf{x},t)=0,
\]
a PINN uses
\[
u_\theta(\mathbf{x},t)\approx u(\mathbf{x},t),
\]
and minimizes
\[
\mathcal{L}(\theta)
=
\mathcal{L}_{\mathrm{data}}
+
\lambda_{\mathrm{PDE}}\mathcal{L}_{\mathrm{PDE}}
+
\lambda_{\mathrm{BC}}\mathcal{L}_{\mathrm{BC}}
+
\lambda_{\mathrm{IC}}\mathcal{L}_{\mathrm{IC}}.
\]

The PDE residual term is
\[
\mathcal{L}_{\mathrm{PDE}}
=
\frac{1}{N_c}\sum_{j=1}^{N_c}
\left|
\mathcal{N}[u_\theta](\mathbf{x}_j,t_j)
\right|^2.
\]
\end{frame}

\begin{frame}{Suggested Equation for Section 1.3.2: Graph Neural Network}
For a mesh graph
\[
\mathcal{G}=(\mathcal{V},\mathcal{E}),
\]
a message-passing GNN can be written as
\[
\mathbf{m}_{ij}^{(\ell)}
=
\phi_e^{(\ell)}
\left(
\mathbf{h}_i^{(\ell)},
\mathbf{h}_j^{(\ell)},
\mathbf{e}_{ij}
\right),
\]
\[
\mathbf{m}_i^{(\ell)}
=
\sum_{j\in \mathcal{N}(i)}
\mathbf{m}_{ij}^{(\ell)},
\]
\[
\mathbf{h}_i^{(\ell+1)}
=
\phi_v^{(\ell)}
\left(
\mathbf{h}_i^{(\ell)},\mathbf{m}_i^{(\ell)}
\right).
\]

This should be connected explicitly to unstructured finite-element meshes.
\end{frame}

%================================================
\section{Chapter 2}
%================================================

\begin{frame}{Chapter 2: Methods}
Chapter 2 is very good and the strongest part of the draft.

\medskip

Its main strengths are:
\begin{itemize}
    \item modular decomposition into EM and thermal surrogates,
    \item graph representation of the computational mesh,
    \item physics-informed losses,
    \item explicit training strategies,
    \item verification through analytical and manufactured solutions.
\end{itemize}

\medskip

The chapter is appropriate for MSc level, but several parts would benefit from more explicit equations and clearer justification.
\end{frame}

\begin{frame}{Chapter 2: Main Improvements}
The most important improvements are:
\begin{enumerate}
    \item explain why the graph representation is the right one,
    \item define node and edge features more explicitly,
    \item write the physics-informed loss in full,
    \item clarify the EM surrogate target variable,
    \item explain training choices and alternatives,
    \item make verification criteria quantitative.
\end{enumerate}
\end{frame}

%================================================
\section{Equations for Chapter 2}
%================================================

\begin{frame}{Suggested Equation for Section 2.1: Surrogate Motivation}
A concise way to express the surrogate-learning goal is:
\[
\mathcal{S}_{\mathrm{FEM}}:
(\Omega,\mathbf{p},\mathbf{J}_{\mathrm{src}})
\mapsto
(\mathbf{A},T),
\]
where \(\mathbf{p}\) denotes material and geometric parameters.

The learned surrogate approximates:
\[
\mathcal{S}_{\theta}
\approx
\mathcal{S}_{\mathrm{FEM}},
\]
but with much smaller evaluation cost:
\[
\mathrm{cost}(\mathcal{S}_{\theta})
\ll
\mathrm{cost}(\mathcal{S}_{\mathrm{FEM}}).
\]
\end{frame}

\begin{frame}{Suggested Equation for Section 2.2: Mesh Graph}
Represent the finite-element mesh as
\[
\mathcal{G}=(\mathcal{V},\mathcal{E}),
\]
where:
\[
i\in\mathcal{V}
\quad\leftrightarrow\quad
\text{mesh node},
\]
\[
(i,j)\in\mathcal{E}
\quad\leftrightarrow\quad
\text{mesh adjacency}.
\]

Node features may include:
\[
\mathbf{h}_i^{(0)}
=
\left[
x_i,y_i,z_i,T_i^n,k_i,\rho_i,c_{p,i},\sigma_i,\mu_i
\right].
\]

Edge features may include:
\[
\mathbf{e}_{ij}
=
\left[
\mathbf{x}_j-\mathbf{x}_i,\,
\norm{\mathbf{x}_j-\mathbf{x}_i}
\right].
\]
\end{frame}

\begin{frame}{Suggested Equation for Section 2.3: Modular Decomposition}
The coupled problem can be decomposed as:
\[
\mathbf{A}^{n}
=
\mathcal{M}_{\theta}^{\mathrm{EM}}
\left(
\mathcal{G},\mathbf{J}_{\mathrm{src}},\sigma(T^n),\mu(T^n)
\right),
\]
\[
Q_{\mathrm{Joule}}^{n}
=
\frac{1}{2}\sigma(T^n)\omega^2|\mathbf{A}^{n}|^2,
\]
\[
T^{n+1}
=
\mathcal{M}_{\phi}^{\mathrm{thermal}}
\left(
\mathcal{G},T^n,Q_{\mathrm{Joule}}^n,k(T^n),\rho,c_p
\right).
\]

This explicitly states the modular surrogate structure.
\end{frame}

\begin{frame}{Suggested Equation for Section 2.4.1: Thermal Time-Step Residual}
Using backward Euler,
\[
R_T^{n+1}
=
\rho c_p
\frac{T_\theta^{n+1}-T^n}{\Delta t}
-
\nabla\cdot(k\nabla T_\theta^{n+1})
-
Q^{n+1}.
\]

The thermal physics-informed loss can be based on
\[
\mathcal{L}_{\mathrm{phys}}^T
=
\frac{1}{N_c}\sum_{i=1}^{N_c}
\left|R_T^{n+1}(\mathbf{x}_i)\right|^2.
\]

This should be explicitly connected to the weak form if the residual is computed in FEM-style form.
\end{frame}

\begin{frame}{Suggested Equation for Section 2.4.2: Thermal Loss}
A complete thermal loss may be:
\[
\mathcal{L}_{T}
=
\mathcal{L}_{\mathrm{data}}^T
+
\lambda_{\mathrm{phys}}\mathcal{L}_{\mathrm{phys}}^T
+
\lambda_{\mathrm{BC}}\mathcal{L}_{\mathrm{BC}}^T.
\]

For supervised FEM data:
\[
\mathcal{L}_{\mathrm{data}}^T
=
\frac{1}{N}\sum_{i=1}^N
\left|
T_{\theta,i}^{n+1}
-
T_{\mathrm{FEM},i}^{n+1}
\right|^2.
\]

For boundary points:
\[
\mathcal{L}_{\mathrm{BC}}^T
=
\frac{1}{N_b}\sum_{i\in\partial\Omega}
\left|
\mathcal{B}[T_\theta](\mathbf{x}_i)-g(\mathbf{x}_i)
\right|^2.
\]
\end{frame}

\begin{frame}{Suggested Equation for Section 2.4.3: Thermal Inputs and Outputs}
State explicitly that the thermal GNN learns a map:
\[
\mathcal{M}_{\theta}^{T}:
\left(
\mathcal{G},
\mathbf{T}^n,
\mathbf{Q}^n,
\mathbf{p}
\right)
\mapsto
\mathbf{T}^{n+1}.
\]

For multi-step rollout:
\[
\mathbf{T}^{n+m}
=
\left(\mathcal{M}_{\theta}^{T}\right)^m
\left(
\mathbf{T}^{n}
\right).
\]

This clarifies whether the model is trained for one-step prediction or long-term simulation.
\end{frame}

\begin{frame}{Suggested Equation for Section 2.5.1: EM Weak Form}
For the frequency-domain MQS equation:
\[
\nabla\times(\nu\nabla\times \mathbf{A})
+
i\omega\sigma\mathbf{A}
=
\mathbf{J}_{\mathrm{src}},
\]
multiply by a test function \(\mathbf{v}\):
\[
\int_\Omega
(\nabla\times \mathbf{v})\cdot
\nu(\nabla\times \mathbf{A})\,d\Omega
+
\int_\Omega
i\omega\sigma \mathbf{v}\cdot \mathbf{A}\,d\Omega
=
\int_\Omega
\mathbf{v}\cdot\mathbf{J}_{\mathrm{src}}\,d\Omega.
\]

This is a useful equation to include in the methods chapter.
\end{frame}

\begin{frame}{Suggested Equation for Section 2.5.2: Nondimensionalization}
Introduce characteristic scales:
\[
\mathbf{x}=L\mathbf{x}',
\qquad
\mathbf{A}=A_0\mathbf{A}',
\qquad
\mathbf{J}=J_0\mathbf{J}'.
\]

Then dimensionless groups can be identified, for example:
\[
\chi
=
\omega\sigma\mu L^2.
\]

This parameter measures the relative importance of eddy-current effects and is related to the magnetic diffusion length or skin-depth scale.
\end{frame}

\begin{frame}{Suggested Equation for Section 2.5.3: Magnetostatic Limit}
In the magnetostatic limit,
\[
\omega\sigma\mathbf{A}\approx 0,
\]
so
\[
\nabla\times(\nu\nabla\times\mathbf{A})
=
\mathbf{J}_{\mathrm{src}}.
\]

This should be clearly identified as the low-frequency or negligible eddy-current limit.
\end{frame}

\begin{frame}{Suggested Equation for Section 2.5.4: Eddy-Current Case}
In the eddy-current case, retain the conductive term:
\[
\nabla\times(\nu\nabla\times\mathbf{A})
+
i\omega\sigma\mathbf{A}
=
\mathbf{J}_{\mathrm{src}}.
\]

The model output may be complex:
\[
\mathbf{A}=\mathbf{A}_R+i\mathbf{A}_I.
\]

Thus the network may predict:
\[
\mathcal{M}_{\theta}^{\mathrm{EM}}:
\mathcal{G}\mapsto
(\mathbf{A}_R,\mathbf{A}_I).
\]

This should be stated explicitly.
\end{frame}

\begin{frame}{Suggested Equation for Section 2.5.6: EM Loss}
A natural EM loss is:
\[
\mathcal{L}_{\mathrm{EM}}
=
\mathcal{L}_{\mathrm{data}}^{A}
+
\lambda_{\mathrm{phys}}\mathcal{L}_{\mathrm{phys}}^{A}
+
\lambda_{\mathrm{BC}}\mathcal{L}_{\mathrm{BC}}^{A}.
\]

For complex-valued vector potential:
\[
\mathcal{L}_{\mathrm{data}}^{A}
=
\frac{1}{N}\sum_i
\left(
\norm{\mathbf{A}_{R,\theta,i}-\mathbf{A}_{R,\mathrm{FEM},i}}^2
+
\norm{\mathbf{A}_{I,\theta,i}-\mathbf{A}_{I,\mathrm{FEM},i}}^2
\right).
\]
\end{frame}

\begin{frame}{Suggested Equation for Section 2.6: Error Metrics}
Use quantitative error metrics such as:
\[
\epsilon_{L^2}
=
\frac{
\norm{\mathbf{u}_\theta-\mathbf{u}_{\mathrm{ref}}}_2
}{
\norm{\mathbf{u}_{\mathrm{ref}}}_2
},
\]
\[
\epsilon_{\infty}
=
\frac{
\norm{\mathbf{u}_\theta-\mathbf{u}_{\mathrm{ref}}}_\infty
}{
\norm{\mathbf{u}_{\mathrm{ref}}}_\infty
}.
\]

For time-dependent rollouts:
\[
\epsilon_{\mathrm{rollout}}(m)
=
\frac{
\norm{\mathbf{T}_\theta^{n+m}-\mathbf{T}_{\mathrm{FEM}}^{n+m}}_2
}{
\norm{\mathbf{T}_{\mathrm{FEM}}^{n+m}}_2
}.
\]
\end{frame}

\begin{frame}{Suggested Equation for Section 2.6.1: Manufactured Solution}
For the method of manufactured solutions, choose an exact function \(u_{\mathrm{ex}}\), then define the forcing:
\[
f(\mathbf{x},t)
=
\mathcal{N}[u_{\mathrm{ex}}](\mathbf{x},t).
\]

Then the numerical or learned solution is compared against \(u_{\mathrm{ex}}\).

For example, for heat conduction:
\[
T_{\mathrm{ex}}(\mathbf{x},t)
=
\sin(\pi x)\sin(\pi y)e^{-t},
\]
and
\[
Q_{\mathrm{MMS}}
=
\rho c_p \frac{\partial T_{\mathrm{ex}}}{\partial t}
-
\nabla\cdot(k\nabla T_{\mathrm{ex}}).
\]
\end{frame}

%================================================
\section{Language and style}
%================================================

\begin{frame}{Language and Style: Strengths}
The writing is generally:
\begin{itemize}
    \item clear,
    \item technically competent,
    \item concise,
    \item appropriate for an MSc thesis.
\end{itemize}

\medskip

The terminology is mostly correct, and the chapter structure is understandable.
\end{frame}

\begin{frame}{Language and Style: Improvements}
Recommended improvements:
\begin{itemize}
    \item avoid overly generic phrases such as ``used in many applications'',
    \item add transitions between theoretical topics,
    \item explain why each equation is relevant to the thesis,
    \item split dense technical paragraphs,
    \item make the central contribution more explicit.
\end{itemize}

\medskip

The writing should move from:
\[
\text{summary of known methods}
\]
toward:
\[
\text{argument for the chosen computational strategy}.
\]
\end{frame}

%================================================
\section{Level assessment}
%================================================

\begin{frame}{Level Assessment}
The level is appropriate for a MSc thesis in computational science.

\medskip

Assessment:
\begin{center}
\begin{tabular}{lc}
\toprule
Category & Assessment \\
\midrule
Physics & Strong MSc level \\
Numerical methods & Strong MSc level \\
Machine learning integration & Above average MSc level \\
Mathematical presentation & Good, but should be expanded \\
Scientific narrative & Good, but too compressed \\
\bottomrule
\end{tabular}
\end{center}
\end{frame}

\begin{frame}{Final Verdict}
The thesis is already a very good MSc-level project.

\medskip

With the improvements suggested here, especially:
\begin{itemize}
    \item stronger mathematical derivations,
    \item clearer graph-learning formulation,
    \item explicit physics-informed losses,
    \item more critical discussion of limitations,
\end{itemize}

it could become a top-tier MSc thesis and approach early PhD-level quality in parts.
\end{frame}

\begin{frame}{Five Highest-Impact Improvements}
The most important revisions are:
\begin{enumerate}
    \item strengthen the narrative flow in Chapter 1,
    \item add one full derivation of a weak form,
    \item write all physics-informed losses explicitly,
    \item expand the graph representation section,
    \item add a critical comparison with FEM and alternative surrogate methods.
\end{enumerate}
\end{frame}

\end{document}
